% HandoutsTexTemplate
% Copyright (C) 2026 Giulio Salvi
%
% This is free software: you can redistribute it and/or modify
% it under the terms of the GNU General Public License as published by
% the Free Software Foundation, either version 3 of the License, or
% (at your option) any later version.

% This software is distributed in the hope that it will be useful,
% but WITHOUT ANY WARRANTY; without even the implied warranty of
% MERCHANTABILITY or FITNESS FOR A PARTICULAR PURPOSE.  See the
% GNU General Public License for more details.
%
% You should have received a copy of the GNU General Public License
% along with this program.  If not, see <https://www.gnu.org/licenses/>.

% Citeproc's CSL interface, following Pandoc's default common.latex partial.
\NewDocumentCommand{\citeproctext}{}{}
\NewDocumentCommand{\citeproc}{mm}{%
  \begingroup\def\citeproctext{#2}\cite{#1}\endgroup
}
\makeatletter
\let\@cite@ofmt\@firstofone
\def\@biblabel#1{}
\def\@cite#1#2{{#1\if@tempswa , #2\fi}}
\makeatother
\newlength{\cslhangindent}
\setlength{\cslhangindent}{1.5em}
\newlength{\csllabelwidth}
\setlength{\csllabelwidth}{3em}
\NewDocumentEnvironment{CSLReferences}{m m}{%
  \begin{list}{}{%
    \setlength{\itemindent}{0pt}
    \setlength{\leftmargin}{0pt}
    \setlength{\parsep}{0pt}
    \ifodd #1
      \setlength{\leftmargin}{\cslhangindent}
      \setlength{\itemindent}{-\cslhangindent}
    \fi
    \setlength{\itemsep}{#2\baselineskip}%
  }%
}{\end{list}}
\NewDocumentCommand{\CSLBlock}{m}{\hfill\break\parbox[t]{\linewidth}{\strut\ignorespaces#1\strut}}
\NewDocumentCommand{\CSLLeftMargin}{m}{\parbox[t]{\csllabelwidth}{\strut#1\strut}}
\NewDocumentCommand{\CSLRightInline}{m}{\parbox[t]{\linewidth-\csllabelwidth}{\strut#1\strut}}
\NewDocumentCommand{\CSLIndent}{m}{\hspace{\cslhangindent}#1}
